% Section 4.3: The exterior algebra as a quotient of $T(V)$ | Part: 1 | Depends on: sec:tensors, sec:metric | Status: done | Original lines: 1042-1112
%!TEX root = ../../main.tex
\subsection{The exterior algebra as a quotient of $T(V)$}\label{sec:extquotient}

The explicit formula \eqref{eq:wedge} hides the simple idea behind the wedge. Starting from the tensor algebra and imposing the
single relation ``a vector wedge itself is zero'' produces the exterior algebra abstractly, for \emph{any} vector space
(in particular for $V$ itself, whose elements are then called multivectors; applied to $V^*$ it gives the algebra of forms).
\begin{definition}[Exterior algebra]
Let $\mathcal I_\wedge\subset T(V)$ be the two-sided ideal generated by $\{v\otimes v: v\in V\}$.
The \emph{exterior algebra} is $\bigwedge V=T(V)/\mathcal I_\wedge$; the induced product is denoted $\wedge$
(the \emph{wedge} or \emph{outer product}), so $v\wedge v=0$ for all $v\in V$. The image of $V^{\otimes k}$ is
$\bigwedge^kV$, the space of \emph{$k$-vectors}; $\bigwedge V=\bigoplus_{k=0}^n\bigwedge^kV$.
\end{definition}

\begin{proposition}[Anticommutativity]\label{prop:anti}
For $u,v\in V$: $u\wedge v=-v\wedge u$. More generally, for $A\in\bigwedge^rV$, $B\in\bigwedge^sV$,
\[
 A\wedge B=(-1)^{rs}\,B\wedge A\qquad(\text{graded commutativity}).
\]
\end{proposition}
\begin{proof}
$0=(u+v)\wedge(u+v)=u\wedge u+u\wedge v+v\wedge u+v\wedge v=u\wedge v+v\wedge u$. For the general statement it suffices
to check blades $v_1\wedge\cdots\wedge v_r$ and $w_1\wedge\cdots\wedge w_s$: moving one $w_j$ past $r$ vectors costs $(-1)^r$,
and there are $s$ of them.
\end{proof}

% REORDERED (F9): the universal property now precedes the basis theorem, whose proof uses it.
\begin{theorem}[Universal property]\label{thm:extuniv}
Alternating $k$-linear maps $V^k\to U$ correspond bijectively to linear maps $\bigwedge^kV\to U$,
via $\beta=\tilde\beta\circ\wedge$. In symbols: $\Alt^k(V;U)\cong\Hom(\bigwedge^kV,U)$.
\end{theorem}

\begin{theorem}[Basis, dimension]\label{thm:extbasis}
If $(e_i)$ is a basis of $V$, then $\{e_{i_1}\wedge\cdots\wedge e_{i_k}: i_1<\cdots<i_k\}$ is a basis of $\bigwedge^kV$. Hence
\[
 \dim\textstyle\bigwedge^kV=\binom nk,\qquad\dim\bigwedge V=2^n,\qquad \bigwedge^kV=0\ (k>n).
\]
Moreover $v_1\wedge\cdots\wedge v_k=0$ if and only if $v_1,\dots,v_k$ are linearly dependent.
\end{theorem}
\begin{proof}[Proof sketch]
Spanning follows from Theorem~\ref{thm:tensorbasis}, Proposition~\ref{prop:anti} (reorder indices) and $e_i\wedge e_i=0$ (kill repeats).
Independence: for each increasing multi-index $I$ the alternating multilinear map $(v_1,\dots,v_k)\mapsto\det\bigl(e^{i_a}(v_b)\bigr)$ induces (by
the universal property, Theorem~\ref{thm:extuniv}) a functional that is $1$ on $e_I$ and $0$ on other basis elements.
\end{proof}

\begin{corollary}[Determinant]\label{cor:det}
If $\dim V=n$ and $a_j=a_j^{\,i}e_i$ then
\[
 a_1\wedge\cdots\wedge a_n=\det\bigl(a_j^{\,i}\bigr)\;e_1\wedge\cdots\wedge e_n .
\]
\end{corollary}
\begin{proof}
Expand by multilinearity: $a_1\wedge\cdots\wedge a_n=\sum_{i_1,\dots,i_n}a_1^{i_1}\cdots a_n^{i_n}\,e_{i_1}\wedge\cdots\wedge e_{i_n}$.
Terms with a repeated index vanish; the others are permutations $\sigma$ and $e_{\sigma(1)}\wedge\cdots\wedge e_{\sigma(n)}=\sgn(\sigma)\,e_1\wedge\cdots\wedge e_n$. The resulting sum is the Leibniz formula for $\det$.
\end{proof}

\begin{remark}[Two models of $\bigwedge^k V$]\label{rem:twomodels}
In characteristic $0$ the quotient model is equivalent to the \emph{subspace} of totally antisymmetric tensors in $V^{\otimes k}$ via
\[
 v_1\wedge\cdots\wedge v_k\;\longmapsto\;\frac1{k!}\sum_{\sigma\in S_k}\sgn(\sigma)\,v_{\sigma(1)}\otimes\cdots\otimes v_{\sigma(k)} .
\]
Also, $\bigwedge^k(V^*)\cong\Alt^k(V;\K)$: \emph{the exterior power of the dual space is the space of alternating $k$-forms on $V$}. This is the model used for differential forms in \S\ref{sec:dform}.
\end{remark}


% REWRITTEN: this remark is now the single place that identifies the quotient construction applied to V^* with the
% algebra of forms; the "second description" in Section sec:wedgetensor refers here instead of repeating it.
\begin{remark}[Same algebra, two normalisations of the embedding]\label{rem:normalisation}
Applied to $V^*$, the quotient $T(V^*)/\langle\alpha\otimes\alpha\rangle$ is the algebra $\bigwedge V^*$ of
Sections~\ref{sec:altforms}--\ref{sec:wedgeprops}: the
class of $\alpha\otimes\beta$ \emph{is} $\alpha\wedge\beta$, and anticommutativity follows from the single rule $\alpha\wedge\alpha=0$ by
polarisation, exactly as shown above. How the quotient is \emph{realised inside} $T^kV^*$ is a matter of convention
(Section~\ref{sec:wedgetensor} and the dictionary of conventions in Section~\ref{sec:conventions}).
\end{remark}
